\section{Related Work}
\label{sec:related}

\paragraph{Test-time compute scaling.} Two axes are established. \emph{Reasoning} scaling (chain-of-thought, tree search, and extended thinking \citep{wei2022chain,yao2023tree,deepseek2025r1,snell2024scaling,kimi2025k15}) spends more tokens before committing. \emph{Sampling} scaling (self-consistency and best-of-$N$ \citep{wang2023selfconsistency,brown2024large}) draws multiple attempts and selects one. Both only rework the same frozen weights and prompt (both are \emph{internal} in the sense of \Cref{sec:framework}, which makes precise how they are bounded). We position \emph{interaction} as a third, external axis that imports information the model does not have.

\paragraph{Self-correction and critique.} Iterative refinement with model-generated feedback (Reflexion \citep{shinn2023reflexion}, ReAct \citep{yao2023react}, Self-Refine \citep{madaan2023selfrefine}, CRITIC \citep{gou2024critic}) is widely used, but its reliability is contested: \citet{huang2024selfcorrect} show that models often cannot self-correct reasoning without an external signal, and \citet{kamoi2025correct} chart when correction does and does not help. Our taxonomy makes the distinction explicit: \emph{ungrounded} critique is subject to the internal ceiling, while refinement against an executed or measured observation is not. Our controls (\Cref{sec:feedback}) then show that it is the instrument's coverage of the defect, not the extra critique pass, that carries the gain.

\paragraph{Grounded and execution feedback.} Tool use and execution feedback (Toolformer \citep{schick2023toolformer}, Gorilla \citep{patil2024gorilla}, Voyager \citep{wang2023voyager}, RLEF \citep{gehring2025rlef}, and, for website generation specifically, WebGen-Agent's multi-level visual/functional feedback \citep{webgen2025}), together with the thinking-vs-doing analysis of \citet{shen2025thinking}, establish that acting on an environment helps. We contribute the information account of \emph{why}, a coverage principle that predicts \emph{when}, and, crucially, the requirement that the \emph{evaluation} be grounded too. We also compare harness architectures (single-agent loops \citep{trankiela2026singleagent}, sample-and-select \citep{ruan2025iad}, proposer--reviewer) under a matched budget.

\paragraph{Model-as-judge reliability.} LLM- and VLM-as-judge evaluation is now standard \citep{zheng2023judging}, and its failure modes (position, verbosity, and self-preference biases, and broader reward-model fragility \citep{casper2023open}) are documented. Our finding is sharper and, to our knowledge, not previously isolated: for artifacts whose quality is a measurable \emph{physical} property (the geometry of a rendered layout), a screenshot judge is not merely biased but \emph{structurally blind} (the defect is dropped by its observation channel before judgment begins), so it both fails to score the defect and, used as a reviewer, fails to fix it. This is why a deterministic instrument is necessary, not merely preferable, to detect interaction scaling on visual modalities.

\paragraph{Practitioner evidence at scale.} Industrial deployment corroborates the gap our framework targets. \citet{hong2026aicoding} reports a controlled study at ByteDance ($3$ frontier coding models $\times$ $3$ agent frameworks $\times$ $100$ runs) on a single real product feature: functional \emph{correctness} exceeded $80\%$ for every pair, yet ``deliverability'' axes (usability, reliability, maintainability, performance) collapsed to $40$--$60\%$ and rose to ${\sim}80\%$ only after harness ``infrastructure'' was added, while correctness barely moved. This is the signature we formalize: internal scaling saturates the easy-to-see axis while the interaction loop carries the rest. Tellingly, the report had to define a multi-axis quality metric before the gap was visible at all (an industrial analogue of our grounded-evaluation requirement), and its grounded fixes (e.g.\ browser-driven validation before commit) are instrument observations in our taxonomy, not ungrounded self-review.

\paragraph{Information theory.} The data-processing inequality \citep{cover2006elements} underlies the formal version of the internal ceiling (\Cref{app:ceiling}): any post-processing of the model's outputs (reasoning, ungrounded review) cannot increase information about the target, whereas an instrument observation introduces a new channel.
